\section{A constructed smooth arithmetic principal open in the same coefficient field}

\begin{theorem}[A finite-type characteristic-zero integer domain has a nonempty smooth principal open]
\label{thm:arithmetic-nonempty-smooth-principal-open}
\lean{EnvelopeProperty.ArithmeticSmoothCoefficientBase.formallySmooth_atBot,EnvelopeProperty.ArithmeticSmoothCoefficientBase.exists_nonzero_smooth_away}
\leanok
Let $R$ be the ACTUAL original integral domain of characteristic zero, of
finite type over $\mathbb Z$. Then there exists an ACTUAL nonzero $r\in R$
such that the original localization $R[1/r]$ is smooth over $\mathbb Z$.
It remains a finite-type integer algebra, a domain, and of characteristic zero.
No smooth-localization witness is an input.
\end{theorem}
\begin{proof}
Since $\mathbb Z$ is Noetherian, the original finite-type $R$ is finitely
presented. Localize at its zero prime. This is the original fraction field
$F$, with the injective map from $R$; its characteristic is zero. The composite
$\mathbb Z\to R\to F$ is ESSENTIALLY of finite type, since the latter map is
a localization. Give $F$ its actual rational-field action. Restricting the
base of this same essentially finite-type algebra along
$\mathbb Z\to\mathbb Q\to F$ shows that $F$ is essentially of finite type
over $\mathbb Q$. This does not assert that $F$ is finite type as an algebra
over either $\mathbb Z$ or $\mathbb Q$.
The native perfect-field theorem therefore makes $F$ formally smooth over
$\mathbb Q$. The map $\mathbb Z\to\mathbb Q$ is a localization and hence
formally smooth. Composition gives formal smoothness of the original
zero-prime localization over $\mathbb Z$. The native openness theorem for
the smooth locus of a finitely presented algebra now produces $r\notin(0)$
with $R[1/r]$ smooth. Localization at one element is finitely presented,
hence the composite is of finite type. As $r\ne0$, this localization embeds
the original domain, so is a domain and retains characteristic zero.
\end{proof}

\begin{definition}[The actual smooth localization inside the same original field]
\label{def:arithmetic-actual-localized-coefficient-ring}
\lean{EnvelopeProperty.ArithmeticSmoothCoefficientBase.awayEmbedding,EnvelopeProperty.ArithmeticSmoothCoefficientBase.awayEmbedding_source,EnvelopeProperty.ArithmeticSmoothCoefficientBase.awayEmbedding_injective,EnvelopeProperty.ArithmeticSmoothCoefficientBase.localizedCoefficientRing,EnvelopeProperty.ArithmeticSmoothCoefficientBase.localizedCoefficientEquiv,EnvelopeProperty.ArithmeticSmoothCoefficientBase.localizedCoefficientEquiv_whole_embedding,EnvelopeProperty.ArithmeticSmoothCoefficientBase.le_localizedCoefficientRing,EnvelopeProperty.ArithmeticSmoothCoefficientBase.localizedCoefficientRing_isLocalization,EnvelopeProperty.ArithmeticSmoothCoefficientBase.exists_smooth_localizedCoefficientRing}
\leanok
\uses{thm:arithmetic-nonempty-smooth-principal-open}
Let $R\subset k$ be the actual finite-type integer coefficient subalgebra
of the SAME original characteristic-zero field. For $r\ne0$ as above, take
the actual localization homomorphism $R[1/r]\to k$ extending the original
inclusion, and let $S\subset k$ be its actual image as an integer subalgebra.
This map is injective; its actual image equivalence
$R[1/r]\simeq_R S$ has the WHOLE original embedding as its composite with
$S\hookrightarrow k$. It makes $S$ an actual localization of $R$ away from
the original $r$. The actual inclusion $R\subset S$ retains all original
polynomial coefficients and all specified finite extras. Internally choosing
$r$ produces such an $S$ smooth and finite type over $\mathbb Z$, a domain
of characteristic zero.
\end{definition}
\begin{proof}
The image of the nonzero original $r$ is a unit of the SAME original field,
so the native localization lift constructs the map, agreeing with the
original inclusion on EVERY element of $R$. The injective original
inclusion and the injective source localization imply injectivity of the
lift by the native localization criterion. Restrict this actual map to its
range, and restrict the range's scalars from $R$ to $\mathbb Z$. The native
injective-map range equivalence gives the claimed whole embedding equality
and localization property. Every original element of $R$ is the image of
its original localization class, so $R\subset S$ literally inside $k$.
Transport the produced smoothness and finite type along the actual range
equivalence; the original field inclusion supplies the domain and
characteristic-zero instances. Apply the preceding theorem to choose the
required nonzero $r$ internally. There is no countable subfield or supplied
smooth-base certificate.
\end{proof}

\begin{theorem}[The same original arithmetic model and whole base-change square survive this localization]
\label{thm:arithmetic-smooth-base-same-model}
\lean{EnvelopeProperty.ArithmeticSmoothCoefficientBase.localizedPresentationModel,EnvelopeProperty.ArithmeticSmoothCoefficientBase.localizedPresentationTensorEquiv,EnvelopeProperty.ArithmeticSmoothCoefficientBase.localizedPresentationTensorEquiv_tmul,EnvelopeProperty.ArithmeticSmoothCoefficientBase.localizedPresentationTensorEquiv_whole_baseChange,EnvelopeProperty.ArithmeticSmoothCoefficientBase.localizedPresentationModel_finitePresentation,EnvelopeProperty.ArithmeticSmoothCoefficientBase.exists_smooth_original_presentation_base}
\leanok
\uses{def:arithmetic-actual-localized-coefficient-ring,thm:arithmetic-same-presentation-model,thm:arithmetic-literal-polynomial-quotient-model}
Retain the ACTUAL coefficient ring $R$, descended relations and literal
model $A_R$ constructed from the original presentation and finite extras.
For the actual $S\subset k$ above, put literally $A_S=S\otimes_R A_R$.
There is a constructed actual equivalence
\[ k\otimes_S A_S\simeq_k A \]
whose composite with the ENTIRE native scalar-extension map
$k\otimes_R A_R\to k\otimes_S(S\otimes_R A_R)$ equals the SAME prior
equivalence to the SAME original $A$ as an equality of whole algebra maps.
For every original $c\in k$, $s\in S$ and polynomial $p\in R[x_j]$, it sends
\[ c\otimes(s\otimes\overline p)\longmapsto
       c\,s\,p(z_j), \]
using the original inclusions into $k$. Thus every original coefficient,
variable and the whole original quotient/evaluation diagram are retained.
The same result specializes to the literal original $k[x_j]/(f_i)$;
the arithmetic model after this base extension is finitely presented over $S$.
The nonzero $r$ and smoothness of $S$ are internally produced from the
original finite equations, not supplied as target assumptions.
\end{theorem}
\begin{proof}
Use the native whole algebra equivalence cancelling the consecutive base
changes $R\to S\to k$, and follow it by the SAME preceding arithmetic
model equivalence. Its native formula sends $c\otimes(s\otimes a)$ to
$(s c)\otimes a$. Applying the preceding whole evaluation formula proves
the displayed readback for every polynomial and every coefficient.
Composing with the inverse native base-change equivalence is exactly the
old whole algebra map; no generator-only matching is used. Finite
presentation survives scalar extension by its native theorem. The literal
quotient case uses the native presentation of the SAME original relation
family. Finally apply the actual smooth principal-open producer to this
same finite-type coefficient subalgebra to obtain $r$ internally.
\end{proof}

\begin{corollary}[The produced arithmetic base genuinely supplies the existing length-two Witt point lift]
\label{cor:arithmetic-produced-smooth-base-witt-point}
\lean{EnvelopeProperty.ArithmeticSmoothCoefficientBase.exists_produced_witt_point}
\leanok
\uses{def:arithmetic-actual-localized-coefficient-ring,thm:arithmetic-smooth-base-same-model}
For every bound $N$, this produced SAME smooth arithmetic base $S\subset k$
has an actual maximal ideal with finite residue field of prime
characteristic $p>N$, together with an ACTUAL homomorphism from $S$ to its
length-two Witt ring reducing to the ENTIRE original quotient map.
The composite with the actual $R\to S$ inclusion retains the original
integer coefficient embedding and the same quotient reduction square.
No family, line-bundle, Frobenius lift or vanishing assertion is made.
\end{corollary}
\begin{proof}
The preceding producer supplies actual smoothness, finite type, domain and
characteristic zero for $S$ internally. Apply the existing actual
large-characteristic Witt point-lift theorem to this SAME $S$ and its
nonzero element $1$. Its native formal-smooth lifting construction gives
the whole reduction equality. Compose that equality with the actual
inclusion from $R$ to obtain the retained original coefficient square.
This is a lift of a point of the arithmetic base, and does not descend or
lift the original geometric family or line bundle.
\end{proof}

This closes only the smooth arithmetic BASE prerequisite. Descent and
construction of the same original geometric family, its projective
embedding and line bundle, and subsequent cohomological or Deligne--Illusie
arguments are separate mathematical producers.
